\section{Математические основы и постановка задачи}
\label{sec:notation}

Настоящий раздел~--- центральная теоретическая часть работы. Здесь фиксируются обозначения, формулируется математическая постановка задачи в терминах переноса меры, излагается аппарат вероятностных потоков и условного флоу матчинга, обсуждается связь с теорией оптимального транспорта и приводится подробное сравнение с диффузионными моделями. В завершение раздела изучаются численные методы интегрирования получающегося ОДУ и обосновывается, почему при линейном интерполянте достаточно одношагового метода Эйлера. Нейронная сеть~$f_\theta$, появляющаяся в формулировках,~--- не более чем инструмент аппроксимации искомого поля скоростей; её конкретная архитектура отложена до раздела~\ref{sec:method}.

\subsection{Обозначения}

Строчные полужирные буквы обозначают векторы: $\mathbf{x}, \mathbf{v} \in \mathbb{R}^d$, заглавные полужирные~--- матрицы. Через~$\pi$ (с нижним индексом) обозначается вероятностная мера на~$(\mathbb{R}^d,\mathcal{B}(\mathbb{R}^d))$; в частности, $\pi_0$ и~$\pi_1$~--- начальная и целевая меры. Их плотности (по мере Лебега) обозначаются $p_0$, $p_1$ соответственно; для общей меры $\pi$ пишется $p$ или $\mathrm{d}\pi/\mathrm{d}\lambda$. Для измеримого отображения~$T:\mathbb{R}^d\to\mathbb{R}^d$ через $T_{\#}\pi$ обозначается \emph{пушфорвард} меры~$\pi$: $(T_{\#}\pi)(A) := \pi(T^{-1}(A))$ для $A\in\mathcal{B}(\mathbb{R}^d)$. Символ $\mathcal{U}[0,1]$~--- равномерное распределение на~$[0,1]$, $\mathcal{N}(\mu,\sigma^2)$~--- нормальное; $\Pi(\pi_0,\pi_1)$~--- множество совместных мер с маргиналями~$\pi_0,\pi_1$. Параметры нейронной сети обозначаются через~$\theta$, $\odot$~--- поэлементное умножение. Норма $\|\cdot\|_2$~--- евклидова на $\mathbb{R}^d$, $W_2(\pi_0,\pi_1)$~--- 2-расстояние Вассерштейна.

Термин \emph{нативный} обозначает КТ без контраста; \emph{контрастный}~--- артериальную фазу. Все срезы имеют размер $512\times512$ пикселей, интенсивности задаются в единицах Хаунсфилда (HU).

\subsection{Постановка задачи}

Пусть $\pi_N$ и $\pi_C$~--- вероятностные меры на $\mathbb{R}^d$ ($d=512^2$), описывающие распределения нативных и контрастных КТ-срезов брюшной полости соответственно. Предполагается наличие парных наблюдений $(\mathbf{x}_N,\mathbf{x}_C)\sim\gamma$, где $\gamma\in\Pi(\pi_N,\pi_C)$~--- совместная мера с заданными маргиналями (на практике~--- эмпирическая мера, образованная парами серий одного пациента после регистрации).

\begin{problem}[Транспортная задача в строгой форме]
\label{prob:strict}
Найти измеримое отображение $T:\mathbb{R}^d\to\mathbb{R}^d$, такое что $T_{\#}\pi_0 = \pi_1$, то есть для каждого $\mathbf{x}_0\sim\pi_0$ образ $T(\mathbf{x}_0)$ распределён согласно~$\pi_1$.
\end{problem}

Уже здесь видны фундаментальные сложности, отделяющие математическую формулировку от прикладной:
\begin{enumerate}
    \item \emph{Существование.} При $d\gg 1$ и многомодальных $\pi_0,\pi_1$ существование единственного измеримого~$T$ не гарантировано без дополнительных предположений (например, абсолютной непрерывности~$\pi_0$ относительно меры Лебега~\cite{brenier1991polar}).
    \item \emph{Однозначность.} Любая композиция~$T\circ S$, где $S$ сохраняет $\pi_0$, также решает задачу~\ref{prob:strict}~--- транспортных отображений бесконечно много, и без дополнительного критерия (например, минимизации кинетической энергии) задача недоопределена.
    \item \emph{Нелинейность.} Подмножество $\{T : T_{\#}\pi_0=\pi_1\}$ \emph{не} образует линейного пространства, что исключает напрямую градиентные методы оптимизации.
\end{enumerate}

Флоу матчинг обходит эти трудности через релаксацию: вместо одного отображения ищется однопараметрическое семейство~$\{\phi_t\}_{t\in[0,1]}$, плавно соединяющее тождественное отображение с искомым переносом. Формально:

\begin{problem}[Динамическая релаксация]
\label{prob:dynamic}
Найти семейство диффеоморфизмов $\phi_t:\mathbb{R}^d\to\mathbb{R}^d$, $t\in[0,1]$, такое что $\phi_0=\mathrm{id}$, $(\phi_1)_{\#}\pi_0 = \pi_1$, и порождающее его поле скоростей $u_t:\mathbb{R}^d\to\mathbb{R}^d$ удовлетворяет некоторому критерию регулярности (например, минимизирует кинетическую энергию~--- см.~\S\ref{sec:ot_simple}).
\end{problem}

Данная работа сосредотачивается на математической природе задачи~\ref{prob:dynamic}, а нейронная сеть~$f_\theta$ выступает лишь как универсальный аппроксиматор поля~$u_t$, выбранный по соображениям выразительности и удобства оптимизации.

\subsection{Теория вероятностных потоков}

\subsubsection{Уравнение непрерывности}

\begin{definition}[Вероятностный поток]
Пусть $u:[0,1]\times\mathbb{R}^d\to\mathbb{R}^d$~--- измеримое поле скоростей, удовлетворяющее локально условию Липшица по пространственной переменной. \emph{Поток}, порождаемый~$u$, есть решение задачи Коши:
\begin{equation}
    \frac{\mathrm{d}}{\mathrm{d}t}\phi_t(\mathbf{x}) = u_t(\phi_t(\mathbf{x})), \quad \phi_0(\mathbf{x}) = \mathbf{x}.
    \label{eq:ode_flow}
\end{equation}
Тогда $\phi_t:\mathbb{R}^d\to\mathbb{R}^d$~--- диффеоморфизм для каждого $t\in[0,1]$.
\end{definition}

Эволюция начального распределения~$\pi_0$ под действием потока описывается семейством мер $\pi_t := (\phi_t)_{\#}\pi_0$. Согласно классическому результату из теории переноса (см.~\cite[гл.~5]{villani2008optimal}), плотность $p_t$ меры $\pi_t$ удовлетворяет \emph{уравнению непрерывности}:
\begin{equation}
    \frac{\partial p_t}{\partial t} + \nabla\!\cdot(p_t\,u_t) = 0.
    \label{eq:continuity}
\end{equation}

Уравнение~\eqref{eq:continuity} является ключевым: оно связывает поле скоростей~$u_t$ с эволюцией вероятностной меры. Если~$u_t$ задано~--- $\pi_t$ однозначно определена начальной мерой. Обратно: задав граничные условия~$\pi_0$ и~$\pi_1$, можно искать поле~$u_t$, обеспечивающее нужный перенос. Уравнение~\eqref{eq:continuity} имеет бесконечное число решений (любое поле, совпадающее с~$u_t$ ``по существу'' на носителе $p_t$, эквивалентно); среди них естественно выбирать ``простейшее'' в некоторой норме.

\subsubsection{Связь с оптимальным транспортом}
\label{sec:ot_simple}

Идея этого раздела простая: среди всех способов перевезти распределение~$\pi_0$ в~$\pi_1$ существует один особенный~--- ``прямой''. Чтобы формализовать, какой именно, нужно решить, что значит ``прямой''.

Представим, что массу одной плотности нужно физически перевезти в другую за единицу времени. Если каждая частица движется со скоростью $\mathbf{v}$, то затрачиваемая ею кинетическая энергия пропорциональна $\|\mathbf{v}\|_2^2$. Суммарная энергия по всем частицам и за всё время:
\begin{equation}
    \mathcal{E}(u) = \int_0^1 \int_{\mathbb{R}^d} \lVert u_t(\mathbf{x})\rVert_2^2\, p_t(\mathbf{x})\,\mathrm{d}\mathbf{x}\,\mathrm{d}t
    = \int_0^1 \mathbb{E}_{\mathbf{x}\sim\pi_t}\bigl[\lVert u_t(\mathbf{x})\rVert_2^2\bigr]\,\mathrm{d}t.
    \label{eq:kinetic_energy}
\end{equation}

Здравый смысл подсказывает: чтобы не тратить лишнюю энергию, частицы должны двигаться по прямым с постоянной скоростью. Это и есть содержание классической теоремы Бенаму--Бренье~\cite{benamou2000computational}: минимум~$\mathcal{E}(u)$ среди всех допустимых потоков равен квадрату расстояния Вассерштейна, $W_2^2(\pi_0,\pi_1)$, а минимизирующий поток~--- прямолинейная интерполяция между парами точек~$(\mathbf{x}_0,\mathbf{x}_1)$.

Чтобы понять, какие именно пары соединять прямыми, нужно зафиксировать совместное распределение $\gamma\in\Pi(\pi_0,\pi_1)$~--- то есть указать, ``куда какая частица едет''. Оптимальный выбор~$\gamma^*$ есть решение задачи Канторовича~\cite{kantorovich1942translocation}:
\[
    \gamma^* = \argmin_{\gamma\in\Pi(\pi_0,\pi_1)}\int \lVert \mathbf{x}_1-\mathbf{x}_0\rVert_2^2\,\mathrm{d}\gamma(\mathbf{x}_0,\mathbf{x}_1).
\]
По теореме Бренье~\cite{brenier1991polar}, при абсолютной непрерывности~$\pi_0$ это $\gamma^*$ сосредоточено на графике градиента выпуклой функции~--- то есть оптимальные пары соединены однозначным детерминированным отображением.

\begin{proposition}[Оптимальность прямолинейных траекторий]
\label{prop:linear_opt}
При фиксированном совместном распределении $\gamma\in\Pi(\pi_0,\pi_1)$ поток
\[
    \phi_t(\mathbf{x}_0) = (1-t)\mathbf{x}_0 + t\mathbf{x}_1, \quad (\mathbf{x}_0,\mathbf{x}_1)\sim\gamma,
\]
минимизирует кинетическую энергию~\eqref{eq:kinetic_energy} среди всех потоков с теми же концевыми точками. Соответствующее условное поле скоростей~$u_t(\cdot\mid\mathbf{x}_1)=\mathbf{x}_1-\mathbf{x}_0$ постоянно по~$t$.
\end{proposition}

В прикладной части работы используется не оптимальный по Канторовичу~$\gamma^*$, а \emph{эмпирическое} совместное распределение~--- пары $(\mathbf{x}_N,\mathbf{x}_C)$ срезов одного пациента после регистрации. Такие пары уже образуют разумное приближение к оптимальному сочленению, так как соответствуют одной и той же анатомии. Поэтому полученные траектории~--- ``почти прямые'', а кривизна выученного поля~$u_t$~--- мала. Это техническое наблюдение в дальнейшем (\S\ref{sec:ode_methods}) превращается в утверждение об одношаговости численного интегрирования.

\subsection{Условный флоу матчинг}
\label{sec:cfm_theory}

\subsubsection{Маргинальное поле скоростей}

Пусть $p_t(\mathbf{x})$~--- плотность $\pi_t$. \emph{Истинное} (маргинальное) поле скоростей, порождающее путь $\{\pi_t\}$, по построению Липмана и соавт.~\cite{lipman2022flow} имеет вид:
\begin{equation}
    u_t(\mathbf{x}) = \frac{\int u_t(\mathbf{x}\mid\mathbf{x}_1)\, p_t(\mathbf{x}\mid\mathbf{x}_1)\, p_1(\mathbf{x}_1)\,\mathrm{d}\mathbf{x}_1}{p_t(\mathbf{x})},
    \label{eq:marginal_field}
\end{equation}
где $u_t(\mathbf{x}\mid\mathbf{x}_1)$~--- условное поле скоростей при фиксированной точке назначения~$\mathbf{x}_1$. Непосредственное вычисление~\eqref{eq:marginal_field} невозможно: интеграл в числителе и плотность~$p_t$ в знаменателе аналитически неизвестны. Однако обучение можно вести \emph{симуляционно-свободно}, используя условные величины:

\begin{theorem}[Симуляционно-свободное обучение~{\cite[Theorem~2]{lipman2022flow}}]
\label{thm:sim_free}
Пусть $f_\theta:\mathbb{R}^d\times[0,1]\to\mathbb{R}^d$~--- произвольная параметрическая функция. Маргинальная потеря флоу матчинга
\begin{equation}
    \mathcal{L}_{\mathrm{FM}}(\theta) = \mathbb{E}_{t\sim\mathcal{U}[0,1],\,\mathbf{x}\sim\pi_t}\bigl[\lVert f_\theta(\mathbf{x},t)-u_t(\mathbf{x})\rVert_2^2\bigr]
    \label{eq:fm_loss_marginal}
\end{equation}
и её условный аналог
\begin{equation}
    \mathcal{L}_{\mathrm{CFM}}(\theta) = \mathbb{E}_{t\sim\mathcal{U}[0,1],\,(\mathbf{x}_0,\mathbf{x}_1)\sim\gamma}\bigl[\lVert f_\theta(\mathbf{x}_t,t)-u_t(\mathbf{x}_t\mid\mathbf{x}_1)\rVert_2^2\bigr]
    \label{eq:cfm_loss}
\end{equation}
имеют одинаковые градиенты по~$\theta$: $\nabla_\theta\mathcal{L}_{\mathrm{FM}}=\nabla_\theta\mathcal{L}_{\mathrm{CFM}}$.
\end{theorem}

\begin{proof}[Набросок]
Раскроем квадрат в~\eqref{eq:fm_loss_marginal}:
\[
    \mathbb{E}\bigl[\lVert f_\theta-u_t\rVert^2\bigr]
    = \mathbb{E}\,\lVert f_\theta\rVert^2 - 2\,\mathbb{E}\langle f_\theta, u_t\rangle + \mathbb{E}\,\lVert u_t\rVert^2.
\]
Третий член не зависит от $\theta$ и обнуляется при дифференцировании. Используя~\eqref{eq:marginal_field}, получим:
\begin{align*}
    \mathbb{E}_{\mathbf{x}\sim\pi_t}\langle f_\theta(\mathbf{x},t), u_t(\mathbf{x})\rangle
    &= \int f_\theta(\mathbf{x},t)^\top\!\left[\int u_t(\mathbf{x}\mid\mathbf{x}_1) p_t(\mathbf{x}\mid\mathbf{x}_1) p_1(\mathbf{x}_1) \mathrm{d}\mathbf{x}_1\right]\mathrm{d}\mathbf{x}\\
    &= \mathbb{E}_{(\mathbf{x},\mathbf{x}_1)\sim p_t(\cdot\mid\mathbf{x}_1)p_1}\langle f_\theta(\mathbf{x},t), u_t(\mathbf{x}\mid\mathbf{x}_1)\rangle.
\end{align*}
Аналогично преобразуется первый член. Разница между~\eqref{eq:fm_loss_marginal} и~\eqref{eq:cfm_loss} оказывается константой по~$\theta$, что и доказывает совпадение градиентов.
\end{proof}

Теорема~\ref{thm:sim_free} обеспечивает обучаемость: для оценки градиента достаточно сэмплировать пары $(\mathbf{x}_0,\mathbf{x}_1)\sim\gamma$ и точки $\mathbf{x}_t$ из явного условного интерполянта, не симулируя траектории.

\subsubsection{Линейный интерполянт}

При выборе линейного пути $(1-t)\mathbf{x}_0+t\mathbf{x}_1$ условное поле скоростей принимает форму:
\begin{equation}
    \mathbf{x}_t = (1-t)\mathbf{x}_0 + t\mathbf{x}_1, \qquad
    u_t(\mathbf{x}_t\mid\mathbf{x}_1) = \mathbf{x}_1-\mathbf{x}_0.
    \label{eq:linear_interp}
\end{equation}

Заметим, что целевая скорость~$\mathbf{x}_1-\mathbf{x}_0$ не зависит от~$t$. Это означает, что нейронная сеть обучается предсказывать стационарное относительно времени поле, тогда как диффузионные модели вынуждены моделировать существенно зависящий от~$t$ шумовой предиктор $\epsilon_\theta(\mathbf{x}_t,t)$ при $t\in\{0,1,\ldots,T\}$. Численные следствия этого факта обсуждаются в~\S\ref{sec:ode_methods}.

\subsection{Сравнение с диффузионными моделями}
\label{sec:ddpm_comparison}

\subsubsection{Прямой и обратный диффузионные процессы}

В диффузионных вероятностных моделях~(DDPM)~\cite{ho2020denoising, sohl2015deep} \emph{прямой процесс} постепенно зашумляет данные, переводя их в стандартную гауссиану. В дискретной форме задаётся марковской цепью
\begin{equation}
    q(\mathbf{x}_t\mid\mathbf{x}_{t-1}) = \mathcal{N}\!\bigl(\sqrt{1-\beta_t}\,\mathbf{x}_{t-1},\;\beta_t\mathbf{I}\bigr), \quad t=1,\ldots,T,
    \label{eq:ddpm_forward_step}
\end{equation}
с расписанием~$\{\beta_t\}_{t=1}^T\subset(0,1)$. Композиция шагов даёт замкнутую формулу:
\begin{equation}
    q(\mathbf{x}_t\mid\mathbf{x}_0) = \mathcal{N}\!\bigl(\sqrt{\bar\alpha_t}\,\mathbf{x}_0,\;(1-\bar\alpha_t)\mathbf{I}\bigr), \quad \bar\alpha_t = \prod_{s=1}^t (1-\beta_s).
    \label{eq:ddpm_forward}
\end{equation}
При $T\to\infty$ и подходящем выборе $\beta_t$ марковская цепь~\eqref{eq:ddpm_forward_step} сходится к стохастическому дифференциальному уравнению~\cite{song2020score}:
\[
    \mathrm{d}\mathbf{x} = -\tfrac{1}{2}\beta(t)\mathbf{x}\,\mathrm{d}t + \sqrt{\beta(t)}\,\mathrm{d}\mathbf{W}_t,
\]
где $\mathbf{W}_t$~--- стандартное винеровское движение. Соответствующее уравнение Фоккера--Планка для плотности $q_t$:
\[
    \partial_t q_t = \nabla\!\cdot\!\bigl(\tfrac{1}{2}\beta(t)\mathbf{x}\,q_t\bigr) + \tfrac{1}{2}\beta(t)\Delta q_t.
\]

\emph{Обратный процесс} в форме SDE имеет вид (Андерсон, 1982; см.~\cite{song2020score}):
\begin{equation}
    \mathrm{d}\mathbf{x} = \bigl[-\tfrac{1}{2}\beta(t)\mathbf{x} - \beta(t)\nabla_{\mathbf{x}}\log q_t(\mathbf{x})\bigr]\mathrm{d}t + \sqrt{\beta(t)}\,\mathrm{d}\bar{\mathbf{W}}_t,
    \label{eq:reverse_sde}
\end{equation}
где $\bar{\mathbf{W}}_t$~--- обратное винеровское движение. Ключевой ингредиент~--- \emph{скоринг-функция} $s(\mathbf{x},t):=\nabla_{\mathbf{x}}\log q_t(\mathbf{x})$, которую и обучают предсказывать нейронной сетью через денойзинг-скоринг-матчинг~\cite{vincent2011connection, hyvarinen2005estimation}:
\begin{equation}
    \mathcal{L}_{\mathrm{DSM}}(\theta) = \mathbb{E}_{t,\mathbf{x}_0,\bepsilon}\bigl[\lVert s_\theta(\mathbf{x}_t,t)+\bepsilon/\sqrt{1-\bar\alpha_t}\rVert_2^2\bigr], \quad \mathbf{x}_t=\sqrt{\bar\alpha_t}\mathbf{x}_0+\sqrt{1-\bar\alpha_t}\bepsilon.
    \label{eq:dsm}
\end{equation}
На практике вместо~$s_\theta$ обычно обучают эквивалентный шумовой предиктор $\bepsilon_\theta(\mathbf{x}_t,t)\approx -\sqrt{1-\bar\alpha_t}\,s(\mathbf{x}_t,t)$.

\subsubsection{Вероятностный ОДУ-поток диффузии}

Сонг и соавт.~\cite{song2020score} показали, что у обратного SDE~\eqref{eq:reverse_sde} имеется эквивалентный по маргиналям детерминированный спутник:
\begin{equation}
    \frac{\mathrm{d}\mathbf{x}}{\mathrm{d}t} = -\tfrac{1}{2}\beta(t)\bigl[\mathbf{x} + \nabla_{\mathbf{x}}\log q_t(\mathbf{x})\bigr],
    \label{eq:diffusion_pf_ode}
\end{equation}
называемый \emph{вероятностным потоковым ОДУ}. Эта формула проясняет глубокую связь диффузии и флоу матчинга: оба метода обучают поле, порождающее перенос меры, но конкретный вид поля и интерполянта различен.

\subsubsection{Структурные различия}

\begin{note}[Структурное различие]
(i)~ОДУ~\eqref{eq:ode_flow} флоу матчинга детерминировано, тогда как обратный диффузионный SDE~\eqref{eq:reverse_sde} стохастичен; (ii)~целевое поле~\eqref{eq:linear_interp} не зависит от $t$, что существенно снижает сложность задачи аппроксимации; (iii)~прямые траектории при линейном интерполянте обладают \emph{нулевой} кривизной, что позволяет численно интегрировать ОДУ за один шаг (Эйлер); диффузионные ОДУ~\eqref{eq:diffusion_pf_ode} имеют существенно нелинейные траектории, что требует многих шагов.
\end{note}

Ускоренные сэмплеры DDIM~\cite{song2020denoising} сокращают число шагов до 50--100 ценой ухудшения качества. В разделе~\ref{sec:ddim_comparison} приводится прямое сравнение: даже DDIM-50 остаётся в~40~раз медленнее и в~5~раз хуже MFM по MAE.

\subsection{Постановка задачи трансляции КТ}

Применим общую теорию из~\S\S\ref{sec:cfm_theory}--\ref{sec:ddpm_comparison} к конкретной задаче трансляции КТ. Положим $\pi_0=\pi_N$, $\pi_1=\pi_C$ (для прямого направления; обратное получается заменой пределов интегрирования). При линейном интерполянте~\eqref{eq:linear_interp} условная потеря~\eqref{eq:cfm_loss} принимает вид
\begin{equation}
    \mathcal{L}_{\mathrm{MFM}}(\theta) =
    \mathbb{E}_{t\sim\mathcal{U}[0,1]}\;
    \mathbb{E}_{(\mathbf{x}_N,\mathbf{x}_C)\sim\gamma}
    \bigl[\lVert f_\theta\!\left((1-t)\mathbf{x}_N+t\mathbf{x}_C,\;t\right) - (\mathbf{x}_C-\mathbf{x}_N)\rVert_2^2\bigr].
    \label{eq:mfm_loss}
\end{equation}

\paragraph{Двунаправленная трансляция.}
Из симметрии линейного интерполянта~\eqref{eq:linear_interp} следует, что единственная обученная сеть~$f_\theta$ реализует трансляцию в обоих направлениях. Сформулируем это утверждение формально.

\begin{proposition}[Двунаправленность MFM]
\label{prop:bidir}
Пусть $f_{\theta^*}$~--- минимизатор условной потери~\eqref{eq:mfm_loss}, в которой выборка пар $(\mathbf{x}_N,\mathbf{x}_C)\sim\gamma$ симметрична относительно перестановки концевых точек: на каждой итерации с равной вероятностью предъявляется как пара $(\mathbf{x}_0,\mathbf{x}_1)=(\mathbf{x}_N,\mathbf{x}_C)$, так и $(\mathbf{x}_C,\mathbf{x}_N)$. Тогда для любого~$\mathbf{x}\in\mathrm{supp}(\pi_N\cup\pi_C)$ выполнено:
\begin{equation}
    \hat{\mathbf{x}}_C = \mathbf{x} + \int_0^1 f_{\theta^*}(\mathbf{x}_t, t)\,\mathrm{d}t, \qquad
    \hat{\mathbf{x}}_N = \mathbf{x} + \int_1^0 f_{\theta^*}(\mathbf{x}_t, t)\,\mathrm{d}t,
    \label{eq:bidir}
\end{equation}
где $\mathbf{x}_t$~--- решение ОДУ~\eqref{eq:ode_flow} с соответствующим начальным условием. Иными словами, единственное поле~$f_{\theta^*}$ реализует оба отображения $\pi_N\to\pi_C$ и $\pi_C\to\pi_N$, и
\begin{equation}
    f_{\theta^*} = \argmin_{f}\bigl[\mathcal{L}_{\mathrm{MFM}}^{N\to C}(f) + \mathcal{L}_{\mathrm{MFM}}^{C\to N}(f)\bigr].
    \label{eq:bidir_obj}
\end{equation}
\end{proposition}

\begin{proof}
Запишем потерю~\eqref{eq:mfm_loss} в обобщённой форме для пары $(\mathbf{x}_0,\mathbf{x}_1)$:
\[
    \mathcal{L}(\theta;\mathbf{x}_0,\mathbf{x}_1) = \mathbb{E}_{t\sim\mathcal{U}[0,1]}\bigl[\lVert f_\theta((1-t)\mathbf{x}_0+t\mathbf{x}_1,\,t)-(\mathbf{x}_1-\mathbf{x}_0)\rVert_2^2\bigr].
\]
Сделаем замену переменной $s=1-t$. Поскольку $\mathrm{d}s=-\mathrm{d}t$ и $s\sim\mathcal{U}[0,1]$ вместе с~$t$, получаем
\begin{align*}
    \mathcal{L}(\theta;\mathbf{x}_0,\mathbf{x}_1)
    &= \mathbb{E}_{s\sim\mathcal{U}[0,1]}\bigl[\lVert f_\theta(s\mathbf{x}_0+(1-s)\mathbf{x}_1,\,1-s)-(\mathbf{x}_1-\mathbf{x}_0)\rVert_2^2\bigr]\\
    &= \mathbb{E}_{s\sim\mathcal{U}[0,1]}\bigl[\lVert -\bigl(-f_\theta(s\mathbf{x}_0+(1-s)\mathbf{x}_1,\,1-s)\bigr)-(\mathbf{x}_1-\mathbf{x}_0)\rVert_2^2\bigr].
\end{align*}
Введём симметризованную сеть $\tilde f_\theta(\mathbf{x},t):=-f_\theta(\mathbf{x},1-t)$ и пару с переставленными концами $(\mathbf{x}_0',\mathbf{x}_1'):=(\mathbf{x}_1,\mathbf{x}_0)$. Тогда последнее выражение равно $\mathcal{L}(\tilde\theta;\mathbf{x}_0',\mathbf{x}_1')$~--- потере с обращённым направлением.

По условию симметрии выборки $\gamma$ относительно перестановки концов,
\[
    \mathbb{E}_{(\mathbf{x}_0,\mathbf{x}_1)\sim\gamma}\,\mathcal{L}(\theta;\mathbf{x}_0,\mathbf{x}_1)
    = \mathbb{E}_{(\mathbf{x}_0',\mathbf{x}_1')\sim\gamma}\,\mathcal{L}(\theta;\mathbf{x}_1',\mathbf{x}_0')
    = \mathbb{E}_{(\mathbf{x}_0',\mathbf{x}_1')\sim\gamma}\,\mathcal{L}(\tilde\theta;\mathbf{x}_0',\mathbf{x}_1').
\]
Следовательно, минимизатор $\theta^*$ обоюдно оптимизирует и прямое, и обратное направления, что доказывает~\eqref{eq:bidir_obj}.

Для проверки~\eqref{eq:bidir} заметим: подставив начальное условие $\mathbf{x}(0)=\mathbf{x}_N$ и интегрируя~\eqref{eq:ode_flow} от $0$ к $1$, получаем $\hat{\mathbf{x}}_C$. Заменой $s=1-t$ и в силу установленной симметрии $f_{\theta^*}(\mathbf{x},t)=-f_{\theta^*}(\mathbf{x},1-t)\circ\text{(перестановка концов)}$, интегрирование от $1$ к $0$ при начальном условии $\mathbf{x}(1)=\mathbf{x}_C$ даёт $\hat{\mathbf{x}}_N$, что и требовалось.
\end{proof}

Конкретно для задачи трансляции КТ это означает:
\begin{itemize}
    \item нативное~$\to$~контрастное: интегрирование~\eqref{eq:ode_flow} от $t=0$ к $t=1$;
    \item контрастное~$\to$~нативное: интегрирование от $t=1$ к $t=0$.
\end{itemize}
Вместо двух специализированных моделей ($f_\theta^{N\to C}$ и $f_\theta^{C\to N}$) требуется одна, что вдвое снижает затраты на обучение и хранение весов. Этот теоретический вывод подтверждается экспериментально: однонаправленный вариант TimeResNet даёт MAE\,=\,4.22\,HU, что не отличается статистически значимо от 4.25\,HU у двунаправленного MFM (таблица~\ref{tab:unidirectional}).

\subsection{Численные методы интегрирования ОДУ}
\label{sec:ode_methods}

При инференсе уравнение~\eqref{eq:ode_flow} решается численно на отрезке $t\in[0,1]$. Зафиксируем равномерную сетку $\{t_k = kh\}_{k=0}^{N}$, $h=1/N$, и обозначим $\hat{\mathbf{x}}_k\approx\phi_{t_k}(\mathbf{x}_0)$. Задача Коши:
\begin{equation}
    \frac{\mathrm{d}\mathbf{x}}{\mathrm{d}t} = f_\theta(\mathbf{x}, t), \quad \mathbf{x}(0) = \mathbf{x}_0.
    \label{eq:ivp}
\end{equation}

Изложение в этом разделе самодостаточно: вводятся понятия локальной и глобальной погрешности, теоремы устойчивости, разбираются методы Эйлера, средней точки, Рунге--Кутты 2-го и 4-го порядков, обсуждаются неявные и адаптивные схемы, а в конце доказывается, что при линейном интерполянте простейший метод Эйлера~--- оптимальный выбор.

\subsubsection{Локальная и глобальная погрешности}

\begin{definition}[Локальная погрешность усечения]
Для одношаговой схемы $\hat{\mathbf{x}}_{k+1} = \hat{\mathbf{x}}_k + h\,\Phi(\hat{\mathbf{x}}_k, t_k, h)$ \emph{локальная погрешность} определяется как
\[
    \tau_k = \mathbf{x}(t_{k+1}) - \mathbf{x}(t_k) - h\,\Phi(\mathbf{x}(t_k), t_k, h),
\]
где $\mathbf{x}(t)$~--- точное решение задачи~\eqref{eq:ivp}.
Схема имеет \emph{порядок} $p$, если $\lVert\tau_k\rVert = O(h^{p+1})$.
\end{definition}

\begin{definition}[Глобальная погрешность]
\emph{Глобальная погрешность} после $N$ шагов равна $e_N = \lVert\hat{\mathbf{x}}_N - \mathbf{x}(1)\rVert$. Для $L$-Липшицева поля скоростей и схемы порядка~$p$ выполняется фундаментальная оценка~\cite{Runge1895, hairer1993solving}:
\begin{equation}
    e_N \leq \frac{C h^p}{L}\bigl(e^{L} - 1\bigr),
    \label{eq:global_error}
\end{equation}
где константа $C$ зависит от $(p{+}1)$-й производной точного решения~$\mathbf{x}(t)$.
\end{definition}

Оценка~\eqref{eq:global_error} классически выводится из дискретной леммы Грёнвалла: при условии~$\lVert\Phi(\mathbf{y},t,h)-\Phi(\mathbf{z},t,h)\rVert\le L\lVert\mathbf{y}-\mathbf{z}\rVert$ и суммировании локальных ошибок $\tau_k$ по индукции получается, что~$e_N \le \tau_{\max}\cdot\sum_{k=0}^{N-1}(1+hL)^k \le (\tau_{\max}/hL)(e^L-1)$.

\subsubsection{Метод Эйлера}

Явный метод Эйлера~\cite{euler1769institutionum}~--- простейшая одношаговая схема. Он получается из формулы Тейлора первого порядка:
\[
    \mathbf{x}(t+h) = \mathbf{x}(t) + h\,\mathbf{x}'(t) + \frac{h^2}{2}\,\mathbf{x}''(t) + O(h^3) = \mathbf{x}(t) + h\,f_\theta(\mathbf{x}(t), t) + O(h^2).
\]
Отбрасывая члены порядка~$O(h^2)$, получаем схему первого порядка:
\begin{equation}
    \hat{\mathbf{x}}_{k+1} = \hat{\mathbf{x}}_k + h\, f_\theta(\hat{\mathbf{x}}_k,\, t_k).
    \label{eq:euler}
\end{equation}
Локальная погрешность $\tau_k = O(h^2)$, глобальная~--- $O(h)$. При~$N=1$ шаг совпадает со всем промежутком~$[0,1]$:
\[
    \hat{\mathbf{x}}_1 = \mathbf{x}_0 + f_\theta(\mathbf{x}_0,\, 0).
\]
Геометрически: метод Эйлера в каждой точке проводит касательную к траектории и движется по ней до следующего узла. Точность определяется тем, насколько ``прямой'' является истинная траектория~--- именно это свойство выполнено для флоу матчинга с линейным интерполянтом.

\subsubsection{Метод средней точки}

Метод средней точки~\cite{10.1145/320831.320840} использует оценку производной в середине шага:
\begin{equation}
    \mathbf{k}_1 = f_\theta(\hat{\mathbf{x}}_k,\, t_k),
    \qquad
    \hat{\mathbf{x}}_{k+1} = \hat{\mathbf{x}}_k + h\,f_\theta\!\left(\hat{\mathbf{x}}_k + \tfrac{h}{2}\mathbf{k}_1,\; t_k + \tfrac{h}{2}\right).
    \label{eq:midpoint}
\end{equation}
Идея метода~--- использовать наклон не в начальной точке шага (как в Эйлере), а в его середине: это компенсирует кривизну траектории в первом порядке. Из разложения Тейлора:
\[
    f_\theta\!\left(\mathbf{x} + \tfrac{h}{2}\mathbf{k}_1,\, t + \tfrac{h}{2}\right)
    = f_\theta + \tfrac{h}{2}\bigl(\partial_t f_\theta + \mathbf{k}_1^\top \nabla_{\mathbf{x}} f_\theta\bigr) + O(h^2)
    = \mathbf{x}'(t+\tfrac{h}{2}) + O(h^2),
\]
поэтому $\tau_k = O(h^3)$, глобальная погрешность~--- $O(h^2)$. Цена~--- два обращения к~$f_\theta$ на шаг.

\subsubsection{Метод Рунге--Кутты 2-го порядка (RK2 / метод Хойна)}

Метод Хойна~\cite{Runge1895} (вариант RK2) строит прогнозирующий шаг методом Эйлера, а затем уточняет по правилу трапеции:
\begin{equation}
    \mathbf{k}_1 = f_\theta(\hat{\mathbf{x}}_k,\, t_k), \qquad
    \mathbf{k}_2 = f_\theta(\hat{\mathbf{x}}_k + h\,\mathbf{k}_1,\; t_k + h), \qquad
    \hat{\mathbf{x}}_{k+1} = \hat{\mathbf{x}}_k + \frac{h}{2}(\mathbf{k}_1 + \mathbf{k}_2).
    \label{eq:rk2}
\end{equation}
Схема~\eqref{eq:rk2} имеет второй порядок точности, требует двух обращений к~$f_\theta$ за шаг. Идея~--- использовать среднее наклонов в начальной и (приближённой) конечной точках шага.

\subsubsection{Метод Рунге--Кутты 4-го порядка (RK4)}

Классический метод Рунге--Кутты 4-го порядка~\cite{kutta1901beitrag} использует четыре вычисления поля скоростей:
\begin{align}
    \mathbf{k}_1 &= f_\theta(\hat{\mathbf{x}}_k,\; t_k), \notag\\
    \mathbf{k}_2 &= f_\theta\!\left(\hat{\mathbf{x}}_k + \tfrac{h}{2}\mathbf{k}_1,\; t_k + \tfrac{h}{2}\right), \notag\\
    \mathbf{k}_3 &= f_\theta\!\left(\hat{\mathbf{x}}_k + \tfrac{h}{2}\mathbf{k}_2,\; t_k + \tfrac{h}{2}\right), \notag\\
    \mathbf{k}_4 &= f_\theta\!\left(\hat{\mathbf{x}}_k + h\,\mathbf{k}_3,\; t_k + h\right), \notag\\
    \hat{\mathbf{x}}_{k+1} &= \hat{\mathbf{x}}_k + \frac{h}{6}\left(\mathbf{k}_1 + 2\mathbf{k}_2 + 2\mathbf{k}_3 + \mathbf{k}_4\right).
    \label{eq:rk4}
\end{align}
Схема имеет четвёртый порядок точности ($\tau_k = O(h^5)$, $e_N = O(h^4)$). Структура весов задаётся таблицей Бутчера:
\[
\begin{array}{c|cccc}
0 & & & & \\
1/2 & 1/2 & & & \\
1/2 & 0 & 1/2 & & \\
1 & 0 & 0 & 1 & \\
\hline
& 1/6 & 1/3 & 1/3 & 1/6
\end{array}
\]
Веса $(1/6, 1/3, 1/3, 1/6)$ соответствуют формуле Симпсона интегрирования $\int_0^1 g(s)\,\mathrm{d}s\approx \tfrac{1}{6}(g(0)+4g(1/2)+g(1))$, отсюда и название ``симпсоновский RK4''.

\subsubsection{Неявные и адаптивные методы (краткий обзор)}

Для жёстких систем ОДУ, в которых явные методы требуют экстремально малого шага из-за условий устойчивости, разработаны неявные схемы. Простейшая~--- \emph{неявный (backward) метод Эйлера}:
\[
    \hat{\mathbf{x}}_{k+1} = \hat{\mathbf{x}}_k + h\,f_\theta(\hat{\mathbf{x}}_{k+1},\,t_{k+1}),
\]
требующая на каждом шаге решения нелинейного уравнения относительно~$\hat{\mathbf{x}}_{k+1}$ (методом Ньютона или фиксированной точки). Неявные схемы $A$-устойчивы: область их устойчивости содержит всю левую полуплоскость. На задачах флоу матчинга, где правая часть~$f_\theta$~--- глубокая нейронная сеть, неявные методы практически неприменимы из-за стоимости итераций Ньютона.

\emph{Адаптивные} методы (Дорманд--Принс, RKF45) подбирают шаг~$h$ автоматически по локальной оценке погрешности, сравнивая предсказания схем разного порядка. Они стандарт для общих задач, но в флоу матчинге не дают выигрыша: целевое поле постоянно по~$t$, и любой шаг~$h$ обеспечивает идентичную точность с фиксированной сеткой.

\subsubsection{Точность при линейном поле скоростей}

\begin{proposition}[Точность Эйлера для MFM]
\label{prop:euler_mfm}
Пусть обученная сеть $f_\theta$ аппроксимирует истинное поле~$u_t$ с равномерной ошибкой $\delta = \sup_{t,\mathbf{x}} \lVert f_\theta(\mathbf{x},t) - u_t(\mathbf{x})\rVert_2$, а $u_t$ является $L$-Липшицевым по~$\mathbf{x}$. Тогда глобальная ошибка метода Эйлера с~$N$ шагами ограничена:
\begin{equation}
    e_N \leq \frac{\delta}{L}\bigl(e^{L}-1\bigr) + \frac{C_2\,h}{2L}\bigl(e^L-1\bigr),
    \label{eq:euler_error}
\end{equation}
где $C_2 = \sup_{t,\mathbf{x}} \lVert\partial_t u_t + (u_t\cdot\nabla)u_t\rVert_2$~--- норма ускорения траектории.
\end{proposition}

\begin{proof}[Эскиз]
Локальная ошибка состоит из двух слагаемых: $\tau_k = h\,\delta + \tfrac{h^2}{2}C_2 + O(h^3)$. Первое~--- из-за несовершенства аппроксимации сети, второе~--- стандартный остаток метода Эйлера. Суммирование с применением леммы Грёнвалла даёт~\eqref{eq:euler_error}.
\end{proof}

При линейном интерполянте~\eqref{eq:linear_interp} условная скорость $u_t(\cdot\mid\mathbf{x}_1) = \mathbf{x}_1 - \mathbf{x}_0$ постоянна по $t$ и $\mathbf{x}_t$, поэтому ускорение траектории равно нулю ($C_2 = 0$ в идеале). Следовательно, второе слагаемое в~\eqref{eq:euler_error} обращается в нуль, и ошибка определяется исключительно качеством аппроксимации сетью~$\delta$. Это даёт строгое математическое обоснование тому, что одношаговый Эйлер достаточен. Более того, методы высшего порядка (RK2, RK4) не дают выигрыша при нулевом~$C_2$, но добавляют несколько вычислений~$f_\theta$, что лишь \emph{накапливает} погрешность аппроксимации. Это согласуется с экспериментальной абляцией в таблице~\ref{tab:solver}: одношаговый Эйлер показывает наилучший MAE и SSIM среди всех численных схем.
